% SOURCE_AUTHORITY: sga5_fr_workpass.tex, lines 10683--10714 (LF-normalized unit).
% SOURCE_SHA256: 791F4EFFC5E02832D5D77ED03518C8156D6F07E4C8238B03545DB93D883FBB28
Para o enunciado do corolário seguinte, convenciona-se pôr, para todo $X$-esquema $h:T\to X$, toda parte fechada $Z$ de $T$ e todo complexo $M$ de $A_X$-módulos,
\[
H^*(T,M)=H^*(T,h^*(M)),\qquad
H^*_Z(T,M)=H^*_Z(T,h^*(M)).
\]

\begin{statement}{Corolário 8.3}
Para todo complexo $M$ de $A_X$-módulos e todo $p\in\mathbb Z$, as aplicações
\[
H^p(X,M)\oplus H^{p-2}(Y',M\otimes\mu^{\otimes -1})
\to
H^p(X',M)\oplus H^{p-2d}(Y,M\otimes\mu^{\otimes -d})
\tag{8.3.1}
\]
e
\[
H_Y^p(X,M)\oplus H^{p-2}(Y',M\otimes\mu^{\otimes -1})
\to
H_{Y'}^p(X',M)\oplus H^{p-2d}(Y,M\otimes\mu^{\otimes -d})
\tag{8.3.2}
\]
definidas pela matriz
\[
\begin{pmatrix}
f^*&j_*\\
0&g_*
\end{pmatrix}
\]
são bijeções.
\end{statement}

\emph{Demonstração.} Denotemos provisoriamente por $u$ o morfismo (8.1.1). Levando-se em conta a fórmula de projeção, as flechas (8.3.1) e (8.3.2) são as induzidas nos objetos de coomologia de grau $p$ pelos morfismos $\R\Gamma(X,u\otimes\operatorname{id}_M)$ e $\R\Gamma_Y(X,u\otimes\operatorname{id}_M)$, respectivamente, de onde se deduz a afirmação.
